% Generated draft: requires mathematical and source review.
\documentclass[11pt]{article}
\usepackage[margin=25mm]{geometry}
\usepackage{fontspec}
\setmainfont{Latin Modern Roman}
\usepackage{amsmath,amssymb,mathtools,unicode-math}
\setmathfont{Latin Modern Math}
\usepackage{enumitem,needspace,xurl,hyperref}
\setlength{\parindent}{0pt}
\setlength{\parskip}{5pt}
\setlength{\emergencystretch}{3em}
\newcommand{\N}{\mathbb{N}}
\newcommand{\Z}{\mathbb{Z}}
\newcommand{\Q}{\mathbb{Q}}
\newcommand{\R}{\mathbb{R}}
\newcommand{\C}{\mathbb{C}}
\newcommand{\F}{\mathbb{F}}
\newcommand{\A}{\mathbb{A}}
\newcommand{\PP}{\mathbb{P}}
\newcommand{\E}{\mathbb{E}}
\newcommand{\eps}{\varepsilon}
\newcommand{\dd}{\,\mathrm{d}}
\begin{document}
\section*{UnsolvedMath ID: 6200022}
\textbf{Source problem code:} catalogue entry
\subsection{Definitions and mathematical statement}
Let $G$ act on the source's spaces $X_1$ and $X_2$. The diagonal action on their product is
\[
 g\cdot(x_1,x_2)=(g\cdot x_1,g\cdot x_2).
\]
The record does not specify the classes of groups or spaces, the metrics, or the hypotheses on either action. In a metric setting, one common convex-core requirement would ask for a nonempty closed convex $G$-invariant subset $C\subseteq X_1\times X_2$ with compact quotient $C/G$. The source does not define convex core or prescribe this particular requirement; its product metric and any minimality condition also need to be supplied.

The first question is whether the diagonal action has such a convex core in the intended sense. The source explicitly includes the special case where $G$ is a surface group and $X_1,X_2$ correspond to different hyperbolic structures. It gives no compactness, marking, or boundary assumptions on that special case.

The second question assumes a space $Z$ with cell-like maps to $X_1$ and $X_2$ and asks whether $Z$ can be chosen as the boundary of the core. In the usual compact-space terminology a cell-like map is a continuous surjection with fibers having the shape of a point. The source does not define $Z$, the map category or equivariance, or whether ``boundary'' means visual boundary at infinity or a topological boundary in the product.

The extracted statement says that the maps target $X_i$ themselves, not $\partial_\infty X_i$. That distinction is left explicit. If the intended $X_i$ are noncompact and the proposed core boundary is compact, a surjective map from that boundary to $X_i$ cannot have the usual compact-space meaning. The exact targets and compactness conventions therefore need source checking; no replacement by boundaries of $X_i$ is silently made.
\subsection{Short English statement}
Does a group acting on two spaces have a convex core for its diagonal product action, including the surface-group case from two hyperbolic structures? If it does, can the common cell-like source space be its boundary? The ambient hypotheses, meaning of core, boundary, and map targets are missing from the extracted record.

The supplied summary calls the question open, mentions the Croke--Kleiner/flat-problem context, and reports no resolution located. It gives no precise reduction or theorem. The OPEN-TRIAGE label and 2005 source are retained with current status marked for review.
\subsection{Sources}
\begin{itemize}
\item Local UnsolvedMath registry, permanent ID 6200022, source code catalogue entry, ``Boundaries of Groups and Kleinian Groups --- Problem 22,'' proposed by Thomas Delzant, 2005. Dataset curation and metadata: Ulam AI and UnsolvedMath Contributors, CC BY 4.0.
\item Kapovich, \emph{Problems on Boundaries of Groups and Kleinian Groups} (2005), Problem 22, PDF page 7, \url{https://www.math.ucdavis.edu/~kapovich/EPR/problems.pdf}; normalized pdf-text extraction, accessed 2026-07-29.
\item The supplied summary mentions the Croke--Kleiner/flat-problem context without a separate primary citation. Source-specific provenance marks current status and statement-reuse terms NEEDS\_REVIEW for the public author-hosted PDF, separately from the dataset licence.
\end{itemize}
\end{document}
